% Bagean: metodhe
% Bab: bukti
% Pérangan: pambuka

\documentclass[../../../include/open-logic-section]{subfiles}

\begin{document}

\olfileid{mth}{prf}{int}

\olsection{Pambuka}

Saka pengalaman sinau logika dhasar, panjenengan mbokmenawa wis
kulina karo !!a{derivation} sistem---mbokmenawa sistem
!!{derivation} deduksi alamiah utawa gaya Fitch, utawa sistem wit
bukti. Panjenengan mbokmenawa éling carané mbuktekaké bab ing sistem
iki: mbuktekaké !!a{formula} utawa nuduhaké yèn sawijining argumèn
valid. Kanggo iku, panjenengan nrapaké aturan-aturan sistem nganti
asil pungkasan kang dikarepaké ketemu. Nalika nalar \emph{babagan}
logika, kita uga mbuktekaké bab-bab tartamtu, nanging lumrahé tanpa
!!a{derivation} sistem. Akèh bukti kang bakal kita rembug ditulis
nganggo basa lumrah, kadhangkala kacampuran simbol, dudu mung basa
logika orde kapisan. Nalika wiwitan nyusun bukti kaya iki,
panjenengan bisa bingung: kepriyé mbuktekaké sawijining bab tanpa
!!a{derivation} sistem? Saka ngendi kudu miwiti? Kepriyé ngerti
yèn buktiné bener?

Sadurungé nyoba mbuktekaké, kudu dingertèni apa iku bukti lan
kepriyé nyusuné. Kaya tegesé jenengé, \emph{bukti} dienggo
nuduhaké yèn sawijining pratelan bener. Bayangna pacelathon:
ana wong takon apa saben bilangan prima saliyané loro iku ganjil.
Wangsulan ``iya'' waé durung cukup; dhèwèké bisa takon
\emph{ngapa}. Ing kéné panjenengan perlu mènèhi bukti.

Ing pacelathon saben dina, wangsulan cekak utawa durung jangkep
mbokmenawa cukup. Nanging ing logika lan matématika kita kepengin
bukti kang ketat: nuduhaké bebeneré sawijining pratelan lan
ngilangi \emph{saben} rasa ragu. Mula saben langkah bukti kudu ana
pambenarané, lan pambenaran mau kudu trep: anggepan kang dienggo
pancen kalebu ing pratelan téoréma, dhéfinisi kang ditrapaké
ditrapaké kanthi bener, lan kesimpulan kang dijupuk manut aturan
inferensi kang sah.

Lumrahé kang dibuktèkaké iku sawijining pratelan. Pratelan kaya iki
nduwèni jeneng warna-warna: proposisi, téoréma, lema, utawa korolari.
Proposisi iku pratelan dhasar kang pantes dibuktèkaké: cukup wigati
kanggo dicathet, senajan ora mesthi jero utawa kerep digunakaké.
Téoréma iku proposisi kang wigati banget. Buktiné asring dipérang
dadi sawatara langkah, lan kadhangkala dijenengi miturut wong kang
pisanan mbuktekaké (upamané Téoréma Cantor lan téoréma
L\"owenheim--Skolem) utawa miturut prakara kang dirembug
(upamané téoréma kasampurnan). Lema iku proposisi utawa téoréma
kang dienggo ing bukti asil kang luwih wigati. Nanging kadhangkala
lema uga wigati dhéwé lan dijenengi miturut panemuné
(upamané Lema Zorn). Korolari iku asil kang gampang katut saka
asil liya.

Pratelan kang bakal dibuktèkaké asring nduwèni anggepan kang
njelasaké jinis obyèk kang dirembug. Pratelan bisa diwiwiti kanthi
``Anggepen $!A$ iku !!a{formula} awujud $!B \lif !C$'' utawa
``Anggepen $\Gamma \Proves !A$'', lan sapanunggalané. Iki
\emph{hipotesis} saka proposisi, téoréma, utawa lema;
panjenengan kena nganggep yèn hipotesis mau bener ing bukti.
Hipotesis matesi pratelan kang dibuktèkaké lan ngenalaké jeneng
obyèk kang dirembug. Upamané, yèn proposisi diwiwiti kanthi
``Anggepen $!A$ iku !!a{formula} awujud $!B \lif !C$'',
panjenengan mung mbuktekaké bab rumus-rumus kang wujudé kaya iku,
yaiku kondisional; $!B \lif !C$ dimangertèni minangka kondisional
sembarang kang bakal dirembug ing bukti.

\end{document}
